\chapter{General Introduction}
\label{chap:introduction}

\section{Context and Motivation}

Video surveillance is widely used in cities, campuses, industrial facilities, and transport infrastructure. In many installations, camera streams still converge on a control room where a small team monitors many screens and must notice rare safety-critical events in real time. Sustained attention across multiple feeds declines quickly, and the number of cameras often makes continuous manual monitoring impossible. As a result, operators frequently review footage after an incident rather than preventing it.

Early intelligent-surveillance systems added motion detection, then object detection and video anomaly detection (\gls{vad}), to centralized servers that ingest every stream. Detection accuracy has improved; modern \gls{yolo} models run in real time and weakly supervised \gls{vad} methods perform well on benchmarks such as UCF-Crime \cite{sultani2018ucfcrime}. Centralization still concentrates bandwidth and compute, creates a single point of failure, and moves identifiable footage away from the camera. It can also give operators many alerts without enough context to act on them.

Two lines of work suggest an alternative. \emph{Multi-agent systems} (\gls{mas}) distribute perception and decision-making across cooperating software agents. They also provide coordination mechanisms such as the Contract Net Protocol \cite{smith1980contract} and auction-based task allocation \cite{parker2002alliance}. \emph{Agentic AI} uses large language models (\gls{llm}s) and vision-language models (\gls{vlm}s) to plan, call tools, and reason over multimodal evidence \cite{sapkota2025ai, yao2023react}. In surveillance, such a layer can explain the evidence produced by the perception system without deciding whether an alert should be raised.

This thesis asks how a surveillance system should organize perception, coordination, alerting, and explanation while keeping privacy and operator oversight in view.

\section{Problem Statement}

The problem addressed in this thesis can be stated as follows:

\begin{quote}
\emph{Design, implement, and evaluate a multi-agent intelligent surveillance system for a semi-closed site, in which autonomous perception agents detect security-relevant events from video and audio at the edge; coordinate with one another to corroborate and actively verify uncertain detections; apply deterministic, auditable alerting policies; and generate grounded natural-language explanations for a human operator---all under privacy-by-design constraints.}
\end{quote}

This problem decomposes into five research questions:

\begin{itemize}
    \item \textbf{RQ1 (Architecture).} Does a distributed, hierarchical \gls{mas} architecture measurably improve responsiveness (time-to-alert) and scalability over a centralized sequential pipeline processing the same sensor streams?
    \item \textbf{RQ2 (Coordination).} Does explicit inter-agent coordination---specifically auction-based verification task allocation inspired by the \gls{cnp}---improve alert precision compared to uncoordinated per-sensor alerting, and at what communication cost?
    \item \textbf{RQ3 (Multimodality).} Does fusing audio events with video events increase detection confidence and reduce false alerts compared to vision-only operation?
    \item \textbf{RQ4 (Agentic explanation).} Can an \gls{llm}-based agentic layer generate incident explanations that are strictly grounded in system evidence---i.e., provably free of fabricated references---without influencing or delaying the safety-critical alert decision?
    \item \textbf{RQ5 (Governance).} Can the entire pipeline operate under privacy-by-design constraints (no raw footage leaves the edge, all exported evidence anonymized, all decisions audited) while remaining useful to the operator?
\end{itemize}

\section{Objectives and Contributions}

The overall objective is to deliver a working, evaluated prototype---named \textbf{AURA-MAS} (Agentic Understanding \& Response Architecture --- Multi-Agent System)---together with a reproducible evaluation methodology. The contributions of this thesis are:

\begin{description}
    \item[C1 --- A hierarchical multi-agent surveillance architecture.] A three-layer design (edge perception, site coordination, governance/operator) with six implemented agent roles plus an operator console, a formal event/alert message schema, and a lightweight \gls{mqtt}/Redis communication substrate (Chapter~\ref{chap:design}).
    \item[C2 --- An auction-based active verification protocol.] A single-round auction mechanism, derived from the \gls{cnp}, in which uncertain ``gray-zone'' incident hypotheses trigger a verification task auctioned among camera agents; the winning agent performs a high-scrutiny re-analysis whose outcome adjusts the fused confidence before the alert decision (Chapters~\ref{chap:design}--\ref{chap:implementation}).
    \item[C3 --- Multimodal spatio-temporal late fusion.] A noisy-OR confidence combination scheme with per-modality reliability weights and cross-modal corroboration bonuses, which guarantees that independent supporting evidence strictly increases incident confidence (Chapter~\ref{chap:design}).
    \item[C4 --- A rule-guarded agentic explanation layer.] An \gls{llm}/\gls{vlm} agent that produces natural-language incident reports \emph{after} the deterministic alert decision, with a mechanical guardrail that verifies every cited evidence identifier against the actual event log and falls back to a deterministic template on any violation (Chapter~\ref{chap:implementation}).
    \item[C5 --- A reproducible evaluation methodology and ablation study.] Scenario replay infrastructure with ground-truth manifests, system-level metrics (event-level F1, mean time-to-alert, false alerts per hour, coordination overhead), and a four-way architectural ablation: centralized baseline, \gls{mas} without coordination, \gls{mas} with rule-based scheduling, and \gls{mas} with auction coordination (Chapter~\ref{chap:evaluation}).
\end{description}

\begin{revisionbox}
The contributions are implementation claims, not blanket performance claims.
The final repeated campaign finds no Holm-corrected significant F1 difference
between the four architectures; in particular, the auction's field-of-view
term was not populated during measurement. Chapter~\ref{chap:evaluation}
and Appendix~\ref{app:corrections} record the resulting qualification.
\end{revisionbox}

\section{Scope and Assumptions}

The thesis targets \emph{semi-closed sites}---campuses, warehouses, industrial facilities---where the set of cameras and microphones is known, zones can be declared (restricted areas, entries), and a single security operator team is responsible. The system performs \emph{event detection and alerting}, not person identification: no facial recognition, biometric matching, or person re-identification (\gls{reid}) across time is performed, by deliberate design choice aligned with the \gls{gdpr} data-minimization principle \cite{gdpr2016} and the prohibitions and high-risk provisions of the EU AI Act \cite{euaiact2024}. Evaluation is performed by replaying recorded multi-sensor scenarios with ground-truth annotations rather than live deployment, which enables controlled, repeatable ablations.

\section{Thesis Organization}

The remainder of this thesis is organized as follows. Chapter~\ref{chap:background_mas} introduces the theoretical foundations of multi-agent systems and agentic AI. Chapter~\ref{chap:background_perception} covers the perception building blocks: object detection, multi-object tracking, video anomaly detection, and audio event detection. Chapter~\ref{chap:sota} surveys the state of the art in distributed and multi-agent surveillance, agentic orchestration frameworks, and the legal/ethical governance landscape, and positions our contribution. Chapter~\ref{chap:design} presents the design of AURA-MAS: architecture, agent roles, message schemas, fusion and coordination mechanisms, and privacy-by-design measures. Chapter~\ref{chap:implementation} details the implementation: technology stack, code organization, and the algorithms of each agent. Chapter~\ref{chap:evaluation} describes the experimental setup, metrics, results of the architectural ablation, and discusses findings and limitations. Chapter~\ref{chap:conclusion} concludes and outlines future work.
